% Diego Sebastián Aliaga Damián
% 2024200481M
% Tema 96: Fintech y financiamiento participativo (crowdfunding) en el Perú: oferta y condiciones
% Estilo del Beamer: rojo, negro y amarillo de la Facultad de Economía (tomados de su logo), con la estructura del
% Beamer de la Unidad I (barra de título oscura, pie en tres partes, portada con campos rotulados y tarjetas).
% Lo cargan beamer.tex (proyección) y beamer_impresion.tex (impresión).
\usepackage[T1]{fontenc}
\usepackage[utf8]{inputenc}
\usepackage[spanish,es-noshorthands,es-nodecimaldot,es-tabla]{babel}
\usepackage{tgheros}
\renewcommand{\familydefault}{\sfdefault}
\usepackage{graphicx,booktabs,tabularx,ragged2e,qrcode,csquotes}
\usepackage{tikz}
\usetikzlibrary{arrows.meta,positioning,calc}
\usepackage[most]{tcolorbox}
\usepackage{newunicodechar}
\newunicodechar{≤}{\ensuremath{\leq}}
\newunicodechar{−}{\ensuremath{-}}

% Cifras del artículo (generadas desde salidas/resultados_04.json) y figuras de 04_analisis.py
% Generado por herramientas/articulo/generar_insumos_articulo.py; no editar a mano.
\newcommand{\cifra}[1]{\ifcsname cifra@#1\endcsname\csname cifra@#1\endcsname\else\textbf{??#1??}\fi}
\expandafter\def\csname cifra@acfPrevia\endcsname{0.84}
\expandafter\def\csname cifra@bootstrap\endcsname{10\,000}
\expandafter\def\csname cifra@coberturaPct\endcsname{90.3}
\expandafter\def\csname cifra@colocacion2022\endcsname{54.5}
\expandafter\def\csname cifra@colocacion2023\endcsname{77.9}
\expandafter\def\csname cifra@colocacion2024\endcsname{83.4}
\expandafter\def\csname cifra@colocacion2025\endcsname{88.5}
\expandafter\def\csname cifra@colocacion2026\endcsname{94.2}
\expandafter\def\csname cifra@colocacionEmp\endcsname{89.5}
\expandafter\def\csname cifra@colocacionPer\endcsname{51.4}
\expandafter\def\csname cifra@columnasPanel\endcsname{35}
\expandafter\def\csname cifra@costoAcumulada\endcsname{0.76}
\expandafter\def\csname cifra@costoMediana\endcsname{5.27}
\expandafter\def\csname cifra@costoTrimestres\endcsname{17}
\expandafter\def\csname cifra@cotaMaxAnual\endcsname{0.22}
\expandafter\def\csname cifra@cotaPctMax\endcsname{1.7}
\expandafter\def\csname cifra@cotaPctMaxDesde2024\endcsname{0.2}
\expandafter\def\csname cifra@crecMensualPrevio\endcsname{8.9}
\expandafter\def\csname cifra@eUnoN\endcsname{8}
\expandafter\def\csname cifra@eUnoPMin\endcsname{0.116}
\expandafter\def\csname cifra@eUnoRhoMax\endcsname{0.01}
\expandafter\def\csname cifra@eUnoRhoMin\endcsname{-0.22}
\expandafter\def\csname cifra@eeffEscala\endcsname{8}
\expandafter\def\csname cifra@eeffEscalaCorregida\endcsname{3}
\expandafter\def\csname cifra@eeffEscalaExcluida\endcsname{5}
\expandafter\def\csname cifra@fechaCorte\endcsname{septiembre de 2026}
\expandafter\def\csname cifra@fechaInicio\endcsname{enero de 2021}
\expandafter\def\csname cifra@filasPanel\endcsname{414}
\expandafter\def\csname cifra@hDosFisherOR\endcsname{2.84}
\expandafter\def\csname cifra@hDosFisherP\endcsname{0.127}
\expandafter\def\csname cifra@hDosP\endcsname{< 0.001}
\expandafter\def\csname cifra@hDosZ\endcsname{14.88}
\expandafter\def\csname cifra@hUnoICinf\endcsname{5.63}
\expandafter\def\csname cifra@hUnoICsup\endcsname{11.27}
\expandafter\def\csname cifra@hUnoMedianaAntes\endcsname{3.32}
\expandafter\def\csname cifra@hUnoMedianaDespues\endcsname{11.37}
\expandafter\def\csname cifra@hUnoNAntes\endcsname{42}
\expandafter\def\csname cifra@hUnoNDespues\endcsname{10}
\expandafter\def\csname cifra@hUnoP\endcsname{< 0.001}
\expandafter\def\csname cifra@hUnoPCota\endcsname{< 0.001}
\expandafter\def\csname cifra@hUnoPSinLider\endcsname{0.124}
\expandafter\def\csname cifra@hUnoU\endcsname{387}
\expandafter\def\csname cifra@hhi2022\endcsname{8\,882}
\expandafter\def\csname cifra@hhi2023\endcsname{9\,770}
\expandafter\def\csname cifra@hhi2024\endcsname{9\,999}
\expandafter\def\csname cifra@hhi2025\endcsname{9\,977}
\expandafter\def\csname cifra@hhi2026\endcsname{9\,321}
\expandafter\def\csname cifra@inicioFPF\endcsname{junio de 2022}
\expandafter\def\csname cifra@interanualAntes\endcsname{169}
\expandafter\def\csname cifra@interanualDespues\endcsname{27}
\expandafter\def\csname cifra@interanualP\endcsname{1.000}
\expandafter\def\csname cifra@lider\endcsname{Inversiones.io}
\expandafter\def\csname cifra@mesesActividadFPF\endcsname{52}
\expandafter\def\csname cifra@mesesDespues\endcsname{10}
\expandafter\def\csname cifra@mesesObtenidos\endcsname{374}
\expandafter\def\csname cifra@mesesSobreTendencia\endcsname{0}
\expandafter\def\csname cifra@monto2022\endcsname{6.46}
\expandafter\def\csname cifra@monto2023\endcsname{16.24}
\expandafter\def\csname cifra@monto2024\endcsname{54.28}
\expandafter\def\csname cifra@monto2025\endcsname{120.28}
\expandafter\def\csname cifra@monto2026\endcsname{104.71}
\expandafter\def\csname cifra@montoMedio2022\endcsname{58.8}
\expandafter\def\csname cifra@montoMedio2023\endcsname{115.2}
\expandafter\def\csname cifra@montoMedio2024\endcsname{215.4}
\expandafter\def\csname cifra@montoMedio2025\endcsname{278.4}
\expandafter\def\csname cifra@montoMedio2026\endcsname{320.2}
\expandafter\def\csname cifra@montoTotal\endcsname{301.97}
\expandafter\def\csname cifra@noObtenidos\endcsname{40}
\expandafter\def\csname cifra@noObtenidosVentana\endcsname{28}
\expandafter\def\csname cifra@participacionLider\endcsname{98.5}
\expandafter\def\csname cifra@pctEmpresarialPrincipal\endcsname{99.5}
\expandafter\def\csname cifra@pctPersonalPrincipal\endcsname{95.3}
\expandafter\def\csname cifra@perLiderFin\endcsname{9}
\expandafter\def\csname cifra@perLiderPub\endcsname{12}
\expandafter\def\csname cifra@placeboN\endcsname{5}
\expandafter\def\csname cifra@placeboPMax\endcsname{< 0.001}
\expandafter\def\csname cifra@plataformasAutorizadas\endcsname{6}
\expandafter\def\csname cifra@plataformasConFinanciamiento\endcsname{3}
\expandafter\def\csname cifra@plataformasVigentes\endcsname{3}
\expandafter\def\csname cifra@principalPersonal\endcsname{Afluenta}
\expandafter\def\csname cifra@proyFin2022\endcsname{110}
\expandafter\def\csname cifra@proyFin2023\endcsname{141}
\expandafter\def\csname cifra@proyFin2024\endcsname{252}
\expandafter\def\csname cifra@proyFin2025\endcsname{432}
\expandafter\def\csname cifra@proyFin2026\endcsname{327}
\expandafter\def\csname cifra@proyFinTotal\endcsname{1\,262}
\expandafter\def\csname cifra@rDosTendencia\endcsname{0.81}
\expandafter\def\csname cifra@saltoTEA\endcsname{19.2}
\expandafter\def\csname cifra@semilla\endcsname{0481}
\expandafter\def\csname cifra@solicitudesDiagnostico\endcsname{26}
\expandafter\def\csname cifra@solicitudesHTTP\endcsname{1\,702}
\expandafter\def\csname cifra@solicitudesHTTPok\endcsname{1\,702}
\expandafter\def\csname cifra@varMonto2023\endcsname{151.3}
\expandafter\def\csname cifra@varMonto2024\endcsname{234.3}
\expandafter\def\csname cifra@varMonto2025\endcsname{121.6}
\expandafter\def\csname cifra@varMonto2026\endcsname{14.7}

\graphicspath{{../../salidas/figuras/}{../posteres_articulos/img/}}
\newcommand{\repositorio}{https://github.com/Diego-Aliaga-Damian/finanzas1-u2-crowdfunding-peru}

% ---------- Colores de la Facultad de Economía (logo: rojo #D82018, negro #181810, borde amarillo) ----------
\definecolor{rojo}{HTML}{C81E16}
\definecolor{rojoclaro}{HTML}{FBE9E7}
\definecolor{negro}{HTML}{1B1B18}
\definecolor{grisoscuro}{HTML}{3A3A36}
\definecolor{grisclaro}{HTML}{F3F3F1}
\definecolor{gris}{HTML}{6B6B66}
\definecolor{amarillo}{HTML}{F2CF00}
\setbeamercolor{normal text}{fg=negro,bg=white}
\setbeamercolor{structure}{fg=rojo}
\setbeamercolor{frametitle}{fg=white,bg=negro}
\setbeamercolor{item}{fg=rojo}
\setbeamercolor{filete rojo}{bg=rojo}
\setbeamercolor{filete amarillo}{bg=amarillo}
\setbeamerfont{frametitle}{size=\large,series=\bfseries}
\setbeamertemplate{navigation symbols}{}
\setbeamertemplate{itemize items}[square]
\setbeamertemplate{enumerate items}[default]
\setlength{\leftmargini}{1.1em}

% Barra de título negra a todo el ancho, con filete rojo y línea amarilla (como el borde del logo)
\setbeamertemplate{frametitle}{%
  \nointerlineskip
  \begin{beamercolorbox}[wd=\paperwidth,ht=0.95cm,dp=0.32cm,leftskip=0.55cm,rightskip=0.55cm]{frametitle}
    \usebeamerfont{frametitle}\insertframetitle
  \end{beamercolorbox}\nointerlineskip
  \begin{beamercolorbox}[wd=\paperwidth,ht=2.6pt,dp=0pt]{filete rojo}\end{beamercolorbox}\nointerlineskip
  \begin{beamercolorbox}[wd=\paperwidth,ht=0.9pt,dp=0pt]{filete amarillo}\end{beamercolorbox}\vspace{0.1cm}}

% Pie en tres partes: autor (negro) | tema (rojo) | número de diapositiva (gris oscuro)
\setbeamertemplate{footline}{%
  \leavevmode\hbox{\tiny\color{white}%
    \colorbox{negro}{\makebox[\dimexpr0.30\paperwidth-2\fboxsep][l]{\hspace{0.3cm}Diego S. Aliaga Damián}}%
    \colorbox{rojo}{\makebox[\dimexpr0.50\paperwidth-2\fboxsep][c]{Financiamiento participativo en el Perú, 2022-2026 · Finanzas I}}%
    \colorbox{grisoscuro}{\makebox[\dimexpr0.20\paperwidth-2\fboxsep][r]{\insertframenumber/\inserttotalframenumber\hspace{0.3cm}}}}}

% Tarjetas: encabezado negro (normal) o rojo (destacada), cuerpo gris claro
\newtcolorbox{tarjeta}[2][]{enhanced, colback=grisclaro, colframe=negro, coltitle=white, fonttitle=\bfseries\footnotesize,
  title={#2}, arc=1.5pt, boxrule=0.5pt, left=4pt, right=4pt, top=2pt, bottom=2pt, toptitle=1pt, bottomtitle=1pt,
  before upper={\footnotesize\RaggedRight}, #1}
\newtcolorbox{destacada}[2][]{enhanced, colback=rojoclaro, colframe=rojo, coltitle=white, fonttitle=\bfseries\footnotesize,
  title={#2}, arc=1.5pt, boxrule=0.5pt, left=4pt, right=4pt, top=2pt, bottom=2pt, toptitle=1pt, bottomtitle=1pt,
  before upper={\footnotesize\RaggedRight}, #1}

\newcommand{\rotulo}[1]{{\color{rojo}\scriptsize\bfseries\MakeUppercase{#1}}}
\newcommand{\dato}[1]{{\color{rojo}\bfseries #1}}
\newcommand{\fuente}[1]{\par\vspace{0.05cm}{\tiny\color{gris}#1}}
\newcommand{\linea}[2]{{\color{rojo}\footnotesize\bfseries #1}\par{\scriptsize #2}\par\vspace{0.08cm}}
\newcommand{\singuiones}{\raggedright\hyphenpenalty=10000\exhyphenpenalty=10000}
